\DocumentMetadata{
  pdfversion=2.0,
  pdfstandard=ua-2,
  lang=en-US,
  tagging=on,
  tagging-setup={math/setup=mathml-SE,tabsorder=structure}
}
\documentclass[11pt,letterpaper]{article}
\usepackage[margin=0.68in,includefoot]{geometry}
\usepackage{newcomputermodern}
\usepackage{amsmath,amscd,mathtools}
\usepackage{tikz,adjustbox}
\providecommand{\square}{\mdlgwhtsquare}
\usepackage[hidelinks]{hyperref}
\hypersetup{
  pdftitle={Functional Analysis PhD Exam, September 2011},
  pdfauthor={University of Florida Department of Mathematics},
  pdfsubject={Graduate mathematics examination},
  pdfdisplaydoctitle=true,
  bookmarksopen=true
}
\setcounter{secnumdepth}{0}
\setlength{\parindent}{0pt}
\setlength{\parskip}{0.28em}
\setlength{\emergencystretch}{3em}
\setlength{\leftmargini}{2.15em}
\setlength{\leftmarginii}{2.65em}
\setlength{\leftmarginiii}{3em}
\pagestyle{empty}
\raggedbottom
\allowdisplaybreaks
\NewTaggingSocketPlug{block/list/label}{examproblem}{%
  \tagstructbegin{tag=Lbl}\tagstructend%
  \tagstructbegin{tag=\LBody}%
  \tagstructbegin{tag=\ProblemHeadingTag,title-o={\ProblemHeadingText}}%
  \tagmcbegin{tag=\ProblemHeadingTag,actualtext={\ProblemHeadingText}}%
  #2%
  \tagmcend%
  \tagstructend%
}
\newenvironment{examproblems}[2]{%
  \def\ProblemHeadingTag{#2}%
  \AssignTaggingSocketPlug{block/list/label}{examproblem}%
  \begin{enumerate}%
  \setcounter{enumi}{#1}%
  \renewcommand{\labelenumi}{\textbf{\theenumi.}}%
  \setlength{\topsep}{0.35em}%
  \setlength{\partopsep}{0pt}%
  \setlength{\itemsep}{0.48em}%
  \setlength{\parsep}{0.18em}%
}{\end{enumerate}}
\newenvironment{examsubparts}{%
  \AssignTaggingSocketPlug{block/list/label}{default}%
  \begin{enumerate}%
  \setlength{\topsep}{0.18em}%
  \setlength{\partopsep}{0pt}%
  \setlength{\itemsep}{0.16em}%
  \setlength{\parsep}{0.1em}%
}{\end{enumerate}}
\newenvironment{examinstructions}{%
  \par\begingroup\itshape
}{%
  \par\endgroup\vspace{0.18em}
}
\makeatletter
\renewcommand\subsection{\@startsection{subsection}{2}{0pt}%
  {-1.25ex plus -.25ex minus -.1ex}%
  {0.55ex plus .1ex}%
  {\normalfont\large\bfseries}}
\renewcommand{\@maketitle}{%
  \newpage\null\vskip -1em%
  {\footnotesize\bfseries
    \href{https://www.ufl.edu/}{University of Florida}, \href{https://math.ufl.edu/}{Department of Mathematics}\par}%
  \vskip -0.5em%
  \tagstructbegin{tag=H1,title={Functional Analysis PhD Exam, September 2011}}%
  \tagpdfparaOff%
  \tagmcbegin{tag=H1}%
    {\LARGE\bfseries\href{https://gma.math.ufl.edu/exams/functional-analysis-phd/}{Functional Analysis PhD Exam}, September 2011\par}%
  \tagmcend%
  \tagpdfparaOn%
  \tagstructend%
  \vskip 0.25em%
  \tagmcbegin{artifact}%
    \hrule height 1pt%
  \tagmcend%
  \vskip 1em%
}
\makeatother
\title{Functional Analysis PhD Exam, September 2011}
\author{}
\date{}
\begin{document}
\maketitle
\thispagestyle{empty}
\begin{examinstructions}
Do any SIX problems. Write solutions in a neat and logical fashion, giving complete reasons for all steps.
\end{examinstructions}
\begin{examproblems}{0}{H2}
\def\ProblemHeadingText{Problem 1.}
\item
Let~\(X\) be a vector space and \(\mathcal F\) a vector space of linear functionals on~\(X\) that separates points.
\begin{examsubparts}
\item[\textnormal{a)}]
Write down a base for the weak topology on~\(X\) induced by \(\mathcal F\).
\item[\textnormal{b)}]
Prove that if a linear functional~\(\varphi\) on~\(X\) is continuous in this topology, then \(\varphi \in \mathcal F\).
\end{examsubparts}
\def\ProblemHeadingText{Problem 2.}
\item
Let~\(X\) be a normed vector space. Suppose \((x_n)\) is a sequence in~\(X\), converging weakly to \(x \in X\). (That is, \(\varphi(x_n) \to \varphi(x)\) for every norm-continuous linear functional~\(\varphi\).) Prove that there is a sequence \(y_n\), where each \(y_n\) is a convex combination of finitely many of the \(x_n\), such that \(\lim \lVert y_n-x\rVert=0\).
\def\ProblemHeadingText{Problem 3.}
\item
Let~\(T\) be a bounded linear operator on a Hilbert space~\(H\). Prove that if both~\(T\) and \(T^*\) are bounded below, then~\(T\) is invertible. (Recall that~\(T\) bounded below means that there is a constant \(c>0\) so that \(\lVert Tx\rVert \geq c\lVert x\rVert\) for all \(x\in H\).)
\def\ProblemHeadingText{Problem 4.}
\item
Let~\(T\) be a bounded linear operator on Hilbert space. Prove that there is a positive operator~\(P\) and a partial isometry~\(U\) such that \(\ker U=\ker T\) and \(T=UP\) (that is, \(T\) has a polar decomposition).
\def\ProblemHeadingText{Problem 5.}
\item
Let~\(T\) be a Hilbert space operator with \(\lVert T\rVert\leq 1\).
\begin{examsubparts}
\item[\textnormal{a)}]
Prove that \(I-T^*T\) and \(I-TT^*\) are positive operators.
\item[\textnormal{b)}]
Prove that
\[
T(I-T^*T)^{1/2}=(I-TT^*)^{1/2}T.
\]
 (Hint: approximate \(\sqrt{1-x}\) by polynomials \(p_n(x)\) and use the spectral theorem.)
\end{examsubparts}
\def\ProblemHeadingText{Problem 6.}
\item
Let \((\lambda_n)\) be a bounded sequence of complex numbers. Consider the weighted shift operator~\(S\) defined on \(\ell^2(\mathbb N)\) by
\[
S(x_0,x_1,x_2,\ldots)=(0,\lambda_0x_0,\lambda_1x_1,\lambda_2x_2,\ldots).
\]
 Characterize those sequences \((\lambda_n)\) for which~\(S\) is compact.
\def\ProblemHeadingText{Problem 7.}
\item
Let~\(A\) be a Banach algebra (over \(\mathbb C\)) and \(a\in A\).
\begin{examsubparts}
\item[\textnormal{a)}]
Define the spectrum of~\(a\).
\item[\textnormal{b)}]
Show that if \(\lambda\in\mathbb C\) and \(|\lambda|>\lVert a\rVert\), then~\(\lambda\) does not lie in the spectrum of~\(a\).
\end{examsubparts}
\end{examproblems}
\end{document}
